\chapter{Materiais e Métodos}
\label{cap:metodos}

\section{Visão Geral da Arquitetura}

A \textit{Rizoma} foi concebida segundo o estilo arquitetural de
microsserviços, com responsabilidades distribuídas entre cinco
componentes principais: o servidor de API (FastAPI), o Worker de
análises (R), o banco de dados transacional (PostgreSQL~16), o motor
de busca analítica (Elasticsearch~8) e o frontend (Next.js~14). A
Figura~\ref{fig:arquitetura} ilustra a topologia geral e o fluxo de
dados entre os componentes.

% \begin{figure}[H]
%   \centering
%   \includegraphics[width=0.9\textwidth]{img/arquitetura.pdf}
%   \caption{Arquitetura geral da \textit{Rizoma}.}
%   \fonte{O autor (2026).}
%   \label{fig:arquitetura}
% \end{figure}

A decisão de não incluir Redis ou outro \textit{broker} de mensagens
externo foi deliberada: a fila de jobs é implementada diretamente no
PostgreSQL via \texttt{LISTEN/NOTIFY}, com \texttt{FOR UPDATE SKIP
LOCKED} garantindo exclusividade de consumo sem concorrência entre
workers. Essa abordagem reduz a superfície operacional e mantém o
PostgreSQL como única fonte da verdade para o estado dos jobs.

\section{Infraestrutura e Ambiente de Implantação}

A plataforma é implantada em um cluster Kubernetes leve (k3s) composto
por cinco nós com processadores AMD. O k3s foi escolhido por sua
pegada reduzida em comparação ao Kubernetes padrão, sem abrir mão da
compatibilidade com a API oficial, essencial para manutenção futura.

Em ambiente de desenvolvimento local, os serviços de backend — PostgreSQL,
Elasticsearch, API e R Worker — são orquestrados via
\texttt{docker-compose}. O frontend Next.js é executado separadamente
com \texttt{npm run dev} na porta 3000.

\subsection{Pipeline de Entrega Contínua}
\label{sec:cicd}

A entrega contínua da plataforma é realizada por um pipeline
\textit{GitHub Actions} que constrói quatro imagens Docker
(\texttt{api}, \texttt{r-worker}, \texttt{postgres}, \texttt{frontend})
e as publica no \textit{GitHub Container Registry} (GHCR). Ao concluir
os builds, um job \texttt{deploy} aciona o webhook do Coolify — uma
plataforma de PaaS auto-hospedada — que puxa as imagens \texttt{:latest}
e recria os contêineres no cluster sem exigir acesso SSH direto ao
servidor de produção. Os domínios de produção são
\texttt{bioapi.flipafile.com} (API) e \texttt{bio.flipafile.com}
(frontend), com TLS gerenciado automaticamente via Let's Encrypt.

A integração entre os dois repositórios (\texttt{bio-platform} e
\texttt{bio-frontend}) é resolvida dentro do workflow: o job
\texttt{frontend} executa uma segunda ação \texttt{actions/checkout}
para clonar \texttt{bio-frontend} antes do \texttt{docker build},
mantendo o contexto de construção inalterado e sem duplicar pipelines
entre repositórios.

\section{Banco de Dados}

\subsection{Modelo de Dados}

O esquema relacional foi construído incrementalmente por meio de cinco
arquivos de migração SQL, aplicados pelo entrypoint do container
PostgreSQL em ordem lexicográfica. As tabelas principais são
descritas no Quadro~\ref{qd:tabelas}.

\begin{table}[H]
\centering
\caption{Principais entidades do modelo relacional.}
\label{qd:tabelas}
\begin{tabularx}{\textwidth}{lX}
\toprule
\textbf{Entidade} & \textbf{Descrição} \\
\midrule
\texttt{projects}         & Projetos científicos (INOVAHERB, Pós-Fogo, Biorremediação), com código único, tipo de marcador (16S ou ITS), status e referência ao usuário criador (\texttt{created\_by}). \\
\texttt{project\_analyses}& Configuração de quais análises e gráficos estão habilitados por projeto. \\
\texttt{samples}          & Amostras FASTQ associadas a projetos, com grupo de tratamento, réplica e OIDs dos Large Objects das leituras R1/R2. \\
\texttt{pipeline\_jobs}   & Fila de jobs com status (\texttt{queued}, \texttt{running}, \texttt{done}, \texttt{failed}), payload JSONB e timestamps. \\
\texttt{analysis\_results}& Resultados analíticos em formato JSONB, produzidos pelo R Worker. \\
\texttt{network\_edges}   & Arestas de redes microbianas com pesos e \textit{keystone score}. \\
\texttt{users}            & Usuários autenticados, com role (\texttt{researcher} ou \texttt{admin}) e URL do avatar Google (\texttt{avatar\_url}), atualizada a cada login. \\
\texttt{invited\_users}   & Convites de acesso com controle de uso (\texttt{used\_at}). \\
\bottomrule
\end{tabularx}
\fonte{O autor (2026).}
\end{table}

\subsection{Fila de Jobs via PostgreSQL}

Um gatilho (\textit{trigger}) na tabela \texttt{pipeline\_jobs} emite
uma notificação via \texttt{pg\_notify('new\_job', id)} sempre que um
novo job é inserido com status \texttt{'queued'}. O R Worker mantém
uma conexão persistente ao banco e aguarda notificações com
\texttt{postgresWaitForNotify(timeout = 30s)}. Ao receber o sinal,
executa \texttt{SELECT ... FOR UPDATE SKIP LOCKED LIMIT 1} para
adquirir o próximo job sem race condition, atualiza o status para
\texttt{'running'} e processa a análise.

\section{API REST}

\subsection{Framework e Organização}

A API foi desenvolvida com FastAPI, framework Python assíncrono
baseado em Starlette e Pydantic~v2. A organização segue o padrão
CQRS: routers de comando (\textit{mutations}) e consulta
(\textit{queries}) são mantidos separados, com a leitura de dados
analíticos servida preferencialmente pelo Elasticsearch, e a escrita
sempre confirmada primeiro no PostgreSQL.

A Tabela~\ref{tab:endpoints} lista os principais endpoints públicos
e autenticados.

\begin{table}[H]
\centering
\caption{Endpoints principais da API v1.}
\label{tab:endpoints}
\begin{tabularx}{\textwidth}{llX}
\toprule
\textbf{Método} & \textbf{Caminho} & \textbf{Descrição} \\
\midrule
GET  & \texttt{/api/v1/projects/}              & Lista projetos com análises configuradas \\
GET  & \texttt{/api/v1/samples/\{project\_id\}} & Lista amostras de um projeto \\
POST & \texttt{/api/v1/samples/upload-pair}     & Upload multipart R1+R2 direto na API; grava como Large Objects \\
POST & \texttt{/api/v1/jobs/enqueue}            & Enfileira análise no pipeline \\
GET  & \texttt{/api/v1/jobs/\{project\_id\}}    & Lista jobs e status \\
WS   & \texttt{/api/v1/jobs/ws/status}          & Feed WebSocket de atualizações em tempo real \\
GET  & \texttt{/api/v1/analysis/\{job\_id\}/results} & Resultados JSON de um job \\
POST & \texttt{/api/v1/samples/import-sra}       & Importa par FASTQ de repositório externo (NCBI SRA ou GEO) \\
GET  & \texttt{/api/v1/samples/sra-preview}     & Retorna metadados de um \textit{accession} sem download \\
GET  & \texttt{/api/v1/samples/fastq-sources}   & Lista repositórios externos disponíveis \\
POST & \texttt{/api/v1/auth/google}             & Troca \textit{access\_token} Google por JWT \\
GET  & \texttt{/api/v1/auth/me}                 & Dados do usuário autenticado \\
\bottomrule
\end{tabularx}
\fonte{O autor (2026).}
\end{table}

\subsection{Importação de Dados Externos — NCBI SRA e GEO}
\label{sec:ncbi}

Para projetos cujas amostras já estão depositadas em repositórios
públicos, a plataforma implementa uma camada de importação baseada no
padrão de \textit{Porta e Adaptador} (Arquitetura Hexagonal)
\cite{cockburn2005hexagonal}. Um protocolo Python (\texttt{FastqSource})
define o contrato que qualquer repositório externo deve cumprir —
dois métodos assíncronos: \texttt{fetch\_metadata} e
\texttt{download\_pair}. Adapters concretos (\texttt{SraAdapter} e
\texttt{GeoAdapter}) implementam esse protocolo para o NCBI SRA/ENA
e para o NCBI GEO, respectivamente, sem que o router da API precise
conhecer qual repositório está sendo acessado.

O \texttt{SraAdapter} consulta a API eUtils do NCBI para obter
metadados do experimento (organismo, \textit{layout} de leitura,
tamanho estimado) e usa a portal ENA para recuperar as URLs de download
dos arquivos FASTQ. O \texttt{GeoAdapter} resolve números de acesso
GSM para identificadores SRR via a relação eLink do NCBI, delegando
então o download ao \texttt{SraAdapter}. Ambos respeitam o limite de
requisições da API pública (0,35 s entre chamadas consecutivas).

Novos repositórios — \textit{ArrayExpress}, bases locais, etc. — podem
ser adicionados ao registro sem alteração do endpoint
\texttt{/import-sra}, pois o registro é carregado de forma
\textit{lazy} em \texttt{app/core/ports.py}.

\subsection{Autenticação}

A autenticação é realizada em duas camadas. No frontend, o NextAuth v5
gerencia o fluxo OAuth2 com o Google, restringindo o acesso ao domínio
\texttt{@ufvjm.edu.br}. O \textit{access\_token} obtido é enviado ao
backend, que o valida consultando o endpoint
\texttt{/oauth2/v3/userinfo} da API do Google e verifica o campo
\texttt{email\_verified}. Em seguida, emite um JWT HS256 com expiração
de 60 minutos contendo os campos \texttt{sub} (UUID do usuário),
\texttt{role} e \texttt{exp}.

O primeiro usuário a fazer login com a tabela \texttt{users} vazia
recebe automaticamente o papel \texttt{admin}. Os demais precisam de
convite prévio cadastrado na tabela \texttt{invited\_users}.

A URL da foto de perfil fornecida pelo Google (\texttt{picture} no
\textit{userinfo}) é persistida em \texttt{users.avatar\_url} a cada
upsert de login. O campo é propagado até o frontend via sessão
NextAuth e exibido no painel lateral do usuário e nos cartões de
projeto, permitindo identificar visualmente o autor de cada projeto.

\section{R Worker}

\subsection{Estrutura do Worker}

O R Worker é um processo R independente que executa um loop de
consumo de mensagens conforme o fluxo apresentado no
Algoritmo~\ref{alg:worker}.

\begin{figure}[H]
\centering
\begin{lstlisting}[language=R, caption={Loop principal do R Worker.}, label={alg:worker}]
repeat {
  notify <- postgresWaitForNotify(conn, timeout = 30)
  if (!is.null(notify)) {
    job <- claim_next_job(conn)   # FOR UPDATE SKIP LOCKED
    if (!is.null(job)) {
      tryCatch(
        process_job(conn, job),
        error = function(e) mark_failed(conn, job$id, e$message)
      )
    }
  }
}
\end{lstlisting}
\fonte{O autor (2026).}
\end{figure}

\subsection{Análises Implementadas}

O dispatcher mapeia o campo \texttt{job\_type} do payload para a
função de análise correspondente. Todos os jobs (exceto FUNGuild em
modo online e GSEA) exigem o campo \texttt{payload\$phyloseq\_key},
que identifica o Large Object \texttt{.rds} no PostgreSQL (campo \texttt{phyloseq\_oid}). O Quadro~\ref{qd:analises}
resume as análises disponíveis.

\begin{table}[H]
\centering
\caption{Análises implementadas no R Worker.}
\label{qd:analises}
\begin{tabularx}{\textwidth}{llX}
\toprule
\textbf{Job type} & \textbf{Marcador} & \textbf{Output} \\
\midrule
\texttt{deseq2}      & 16S/ITS & Tabela de DEGs → \texttt{analysis\_results} \\
\texttt{ancombc2}    & ITS     & Táxons diferenciais → \texttt{analysis\_results} \\
\texttt{maaslin2}    & 16S     & Coeficientes de associação → \texttt{analysis\_results} \\
\texttt{spieceasi}   & 16S/ITS & Rede nós+arestas → \texttt{analysis\_results} + \texttt{network\_edges} \\
\texttt{random\_forest} & 16S/ITS & Modelo \texttt{.rds} → PG Large Object (\texttt{result\_oid}); importância → \texttt{analysis\_results} \\
\texttt{gsea}        & 16S/ITS & Pathways GO/KEGG → \texttt{analysis\_results}; suporte a 8 organismos via \texttt{ORG\_DB\_MAP} e \texttt{KEGG\_ORG\_MAP} \\
\texttt{funguild}    & ITS     & Anotações de guilda → \texttt{analysis\_results} \\
\texttt{picrust2}    & 16S     & Predição de vias metabólicas → \texttt{analysis\_results} \\
\bottomrule
\end{tabularx}
\fonte{O autor (2026).}
\end{table}

Uma decisão arquitetural fundamental é que o R Worker \textbf{nunca
gera figuras}: todo output é dados estruturados (CSV transformado em
JSONB ou JSON de redes). As visualizações são geradas exclusivamente
no frontend via Plotly.js e Cytoscape.js, o que facilita a
customização interativa pelo usuário e elimina dependências de pacotes
gráficos (\texttt{ggplot2}, \texttt{Cairo}) no ambiente R.

\section{Frontend}

\subsection{Stack e Roteamento}

O frontend é construído com Next.js~14 utilizando o App Router e
Server Components onde aplicável. A biblioteca SWR provê cache de
dados com revalidação automática. As principais rotas estão listadas
no Quadro~\ref{qd:rotas}.

\begin{table}[H]
\centering
\caption{Rotas do frontend.}
\label{qd:rotas}
\begin{tabularx}{\textwidth}{lX}
\toprule
\textbf{Rota} & \textbf{Descrição} \\
\midrule
\texttt{/}                    & Dashboard com cards de projetos e métricas \\
\texttt{/projects/[id]}       & Upload FASTQ, upload \texttt{.rds}, enfileiramento de jobs \\
\texttt{/jobs}                & Feed em tempo real via WebSocket \\
\texttt{/diversity}           & PCoA e diversidade alfa (Plotly.js) \\
\texttt{/cross-project}       & Painel com 6 PCoAs (Plotly subplots) \\
\texttt{/analysis/[id]}       & Resultados detalhados por job \\
\texttt{/network/[id]}        & Rede microbiana interativa (Cytoscape.js) \\
\texttt{/admin/users}         & Gestão de usuários e convites \\
\texttt{/admin/projects/new}  & Formulário de criação de projeto (admin) \\
\texttt{/docs}                & Documentação interativa da API (Swagger UI) \\
\bottomrule
\end{tabularx}
\fonte{O autor (2026).}
\end{table}

\subsection{Documentação Interativa da API}

A rota \texttt{/docs} do frontend incorpora o
\textit{Swagger UI}\footnote{swagger-ui-react v5, renderizado inteiramente
no cliente via \texttt{next/dynamic} com \texttt{ssr: false}.},
que consome o endpoint \texttt{/openapi.json} servido nativamente
pelo FastAPI. Isso disponibiliza a documentação dos endpoints —
parâmetros, schemas Pydantic, códigos de resposta — diretamente na
interface da plataforma, sem necessidade de acesso externo ao servidor
de API.

\subsection{Visualizações}

Os dados retornados pela API são renderizados por componentes
Plotly.js (gráficos 2D/3D: PCoA, volcano plot, MA plot, heatmap,
\textit{dot plot}, barras de importância) e Cytoscape.js (rede
microbiana). A página \texttt{/cross-project} utiliza subplots
Plotly para compor o painel de 6 PCoAs que constitui a figura central
do TCC.

\section{Armazenamento de Arquivos — PostgreSQL Large Objects}

Os arquivos binários da plataforma — FASTQs comprimidos, objetos
\texttt{.rds} do phyloseq e relatórios de QC — são armazenados
diretamente no PostgreSQL 16 por meio do mecanismo de \textit{Large
Objects} (LO). Cada arquivo é inserido via \texttt{lo\_import} ou
pela interface de \textit{streaming} do \texttt{asyncpg}, gerando um
OID referenciado nas tabelas \texttt{samples} e
\texttt{pipeline\_jobs}. O Quadro~\ref{qd:lo} descreve a organização
lógica dos dados binários.

\begin{table}[H]
\centering
\caption{Categorias de dados binários armazenados como Large Objects.}
\label{qd:lo}
\begin{tabularx}{\textwidth}{lX}
\toprule
\textbf{Categoria} & \textbf{Descrição} \\
\midrule
FASTQs brutos (R1/R2) & Leituras brutas do sequenciador, referenciadas em \texttt{samples.fastq\_r1\_oid} e \texttt{fastq\_r2\_oid} \\
Phyloseq \texttt{.rds}  & Objeto R serializado após pipeline QIIME2, referenciado em \texttt{pipeline\_jobs.payload} \\
Relatórios QC           & FastQC/MultiQC em HTML, vinculados ao projeto \\
Bases de referência     & SILVA 138 e UNITE v10 para classificação taxonômica \\
\bottomrule
\end{tabularx}
\fonte{O autor (2026).}
\end{table}

Essa abordagem centraliza backup, replicação e controle de acesso
em um único serviço, simplificando a operação do cluster. O acesso
transacional garante que um arquivo nunca fique visível ao R Worker
antes da confirmação do registro correspondente, evitando condições
de corrida inerentes ao uso de storage externo.

\section{Projetos Científicos Suportados}

\subsection{INOVAHERB}

Projeto de análise fatorial do micobioma (marcador ITS) de plantas
medicinais. As amostras seguem a nomenclatura Illumina no formato
\texttt{ITSA\_<num>\_<tratamento>\_L001\_R1\_001.fastq.gz}, onde o
grupo de tratamento segue o padrão \texttt{T\{n\}B\{n\}}. A análise
principal é ANCOM-BC2, com SpiecEasi para inferência de rede.

\subsection{Pós-Fogo}

Monitoramento longitudinal da microbiota bacteriana (16S) em solos
após eventos de queimada. O design experimental é de série temporal,
com coletas em múltiplos pontos pós-fogo. MaAsLin2 é empregado para
modelar as associações entre abundância de táxons, tempo e variáveis
edáficas. PICRUSt2 é usado para predição funcional.

\subsection{Biorremediação}

Correlação de comunidades fúngicas e bacterianas (ITS + 16S) com
concentrações de metais pesados em solos contaminados. FUNGuild
anota as guildas ecológicas dos fungos identificados; Random Forest
seleciona os táxons preditores mais relevantes; ANCOM-BC2 quantifica
as abundâncias diferenciais entre áreas contaminadas e controle.
